% ===========================================================
\section{Dataset}
% ===========================================================

\subsection{Material Selection}

We compiled a dataset of 358 MAX-phase compounds of the
M\textsubscript{2}AX stoichiometry,
covering transition metals M $\in$ \{Ti, V, Cr, Zr, Nb, Mo,
Hf, Ta, W, Sc, Y\},
A-group elements A $\in$ \{Al, Si, P, S, Ga, Ge, As, In, Sn,
Sb, Tl, Pb, Bi, Au, Ag, Cu, Zn, Cd, Hg\},
and X $\in$ \{C, N\}.

\subsection{DFT Calculations}

Crystal structures were optimised at the GGA-PBE level
using VASP~\cite{Kresse1996}.
Interatomic force constants were computed via the
finite-displacement method as implemented in
Phonopy~\cite{Togo2015},
using $2\times2\times1$ supercells.
Phonon dispersion curves were evaluated along the
standard hexagonal high-symmetry path
$\Gamma$--M--K--$\Gamma$--A--L--H--A with 357 $\mathbf{q}$-points.

\subsection{Elastic Constants}

Single-crystal elastic constants ($C_{11}$, $C_{12}$, $C_{13}$,
$C_{33}$, $C_{44}$, $C_{66}$) and derived quantities
(Hill bulk modulus $B$, shear modulus $G$,
Young's modulus $E$, Debye temperature $\Theta_D$)
were computed for all 358 materials using the strain-energy approach.

\subsection{Atomic Feature Vector}

Each atom is represented by a 100-dimensional feature vector
drawn from the periodic-table database \texttt{ptable2}.
Features include electronegativity, covalent radii, valence
electron counts, ionisation energies, and pseudopotential
parameters.
A systematic ablation study (Section~\ref{sec:ablation})
identifies the features most predictive of phonon accuracy.

\subsection{Data Splits}

The 358 structures are split into
train / validation / test sets of
286 / 35 / 37 samples (80/10/10\%)
using a fixed random seed to ensure reproducibility.
